\IfFileExists{stacks-project.cls}{%
\documentclass{stacks-project}
}{%
\documentclass{amsart}
}

% For dealing with references we use the comment environment
\usepackage{verbatim}
\newenvironment{reference}{\comment}{\endcomment}
%\newenvironment{reference}{}{}
\newenvironment{slogan}{\comment}{\endcomment}
\newenvironment{history}{\comment}{\endcomment}

% For commutative diagrams we use Xy-pic
\usepackage[all]{xy}

% We use 2cell for 2-commutative diagrams.
\xyoption{2cell}
\UseAllTwocells

% We use multicol for the list of chapters between chapters
\usepackage{multicol}

% This is generally recommended for better output
\usepackage{lmodern}
\usepackage[T1]{fontenc}

% For cross-file-references
\usepackage{xr-hyper}

% Package for hypertext links:
\usepackage{hyperref}

% For any local file, say "hello.tex" you want to link to please
% use \externaldocument[hello-]{hello}
\externaldocument[introduction-]{introduction}
\externaldocument[conventions-]{conventions}
\externaldocument[sets-]{sets}
\externaldocument[categories-]{categories}
\externaldocument[topology-]{topology}
\externaldocument[sheaves-]{sheaves}
\externaldocument[sites-]{sites}
\externaldocument[stacks-]{stacks}
\externaldocument[fields-]{fields}
\externaldocument[algebra-]{algebra}
\externaldocument[brauer-]{brauer}
\externaldocument[homology-]{homology}
\externaldocument[derived-]{derived}
\externaldocument[simplicial-]{simplicial}
\externaldocument[more-algebra-]{more-algebra}
\externaldocument[smoothing-]{smoothing}
\externaldocument[modules-]{modules}
\externaldocument[sites-modules-]{sites-modules}
\externaldocument[injectives-]{injectives}
\externaldocument[cohomology-]{cohomology}
\externaldocument[sites-cohomology-]{sites-cohomology}
\externaldocument[dga-]{dga}
\externaldocument[dpa-]{dpa}
\externaldocument[sdga-]{sdga}
\externaldocument[hypercovering-]{hypercovering}
\externaldocument[schemes-]{schemes}
\externaldocument[constructions-]{constructions}
\externaldocument[properties-]{properties}
\externaldocument[morphisms-]{morphisms}
\externaldocument[coherent-]{coherent}
\externaldocument[divisors-]{divisors}
\externaldocument[limits-]{limits}
\externaldocument[varieties-]{varieties}
\externaldocument[topologies-]{topologies}
\externaldocument[descent-]{descent}
\externaldocument[perfect-]{perfect}
\externaldocument[more-morphisms-]{more-morphisms}
\externaldocument[flat-]{flat}
\externaldocument[groupoids-]{groupoids}
\externaldocument[more-groupoids-]{more-groupoids}
\externaldocument[etale-]{etale}
\externaldocument[chow-]{chow}
\externaldocument[intersection-]{intersection}
\externaldocument[pic-]{pic}
\externaldocument[weil-]{weil}
\externaldocument[adequate-]{adequate}
\externaldocument[dualizing-]{dualizing}
\externaldocument[duality-]{duality}
\externaldocument[discriminant-]{discriminant}
\externaldocument[derham-]{derham}
\externaldocument[local-cohomology-]{local-cohomology}
\externaldocument[algebraization-]{algebraization}
\externaldocument[curves-]{curves}
\externaldocument[resolve-]{resolve}
\externaldocument[models-]{models}
\externaldocument[functors-]{functors}
\externaldocument[equiv-]{equiv}
\externaldocument[pione-]{pione}
\externaldocument[etale-cohomology-]{etale-cohomology}
\externaldocument[proetale-]{proetale}
\externaldocument[relative-cycles-]{relative-cycles}
\externaldocument[more-etale-]{more-etale}
\externaldocument[trace-]{trace}
\externaldocument[crystalline-]{crystalline}
\externaldocument[spaces-]{spaces}
\externaldocument[spaces-properties-]{spaces-properties}
\externaldocument[spaces-morphisms-]{spaces-morphisms}
\externaldocument[decent-spaces-]{decent-spaces}
\externaldocument[spaces-cohomology-]{spaces-cohomology}
\externaldocument[spaces-limits-]{spaces-limits}
\externaldocument[spaces-divisors-]{spaces-divisors}
\externaldocument[spaces-over-fields-]{spaces-over-fields}
\externaldocument[spaces-topologies-]{spaces-topologies}
\externaldocument[spaces-descent-]{spaces-descent}
\externaldocument[spaces-perfect-]{spaces-perfect}
\externaldocument[spaces-more-morphisms-]{spaces-more-morphisms}
\externaldocument[spaces-flat-]{spaces-flat}
\externaldocument[spaces-groupoids-]{spaces-groupoids}
\externaldocument[spaces-more-groupoids-]{spaces-more-groupoids}
\externaldocument[bootstrap-]{bootstrap}
\externaldocument[spaces-pushouts-]{spaces-pushouts}
\externaldocument[spaces-chow-]{spaces-chow}
\externaldocument[groupoids-quotients-]{groupoids-quotients}
\externaldocument[spaces-more-cohomology-]{spaces-more-cohomology}
\externaldocument[spaces-simplicial-]{spaces-simplicial}
\externaldocument[spaces-duality-]{spaces-duality}
\externaldocument[formal-spaces-]{formal-spaces}
\externaldocument[restricted-]{restricted}
\externaldocument[spaces-resolve-]{spaces-resolve}
\externaldocument[formal-defos-]{formal-defos}
\externaldocument[defos-]{defos}
\externaldocument[cotangent-]{cotangent}
\externaldocument[examples-defos-]{examples-defos}
\externaldocument[algebraic-]{algebraic}
\externaldocument[examples-stacks-]{examples-stacks}
\externaldocument[stacks-sheaves-]{stacks-sheaves}
\externaldocument[criteria-]{criteria}
\externaldocument[artin-]{artin}
\externaldocument[quot-]{quot}
\externaldocument[stacks-properties-]{stacks-properties}
\externaldocument[stacks-morphisms-]{stacks-morphisms}
\externaldocument[stacks-limits-]{stacks-limits}
\externaldocument[stacks-cohomology-]{stacks-cohomology}
\externaldocument[stacks-perfect-]{stacks-perfect}
\externaldocument[stacks-introduction-]{stacks-introduction}
\externaldocument[stacks-more-morphisms-]{stacks-more-morphisms}
\externaldocument[stacks-geometry-]{stacks-geometry}
\externaldocument[moduli-]{moduli}
\externaldocument[moduli-curves-]{moduli-curves}
\externaldocument[examples-]{examples}
\externaldocument[exercises-]{exercises}
\externaldocument[guide-]{guide}
\externaldocument[desirables-]{desirables}
\externaldocument[coding-]{coding}
\externaldocument[obsolete-]{obsolete}
\externaldocument[fdl-]{fdl}
\externaldocument[index-]{index}

% Theorem environments.
%
\theoremstyle{plain}
\newtheorem{theorem}[subsection]{Theorem}
\newtheorem{proposition}[subsection]{Proposition}
\newtheorem{lemma}[subsection]{Lemma}

\theoremstyle{definition}
\newtheorem{definition}[subsection]{Definition}
\newtheorem{example}[subsection]{Example}
\newtheorem{exercise}[subsection]{Exercise}
\newtheorem{situation}[subsection]{Situation}

\theoremstyle{remark}
\newtheorem{remark}[subsection]{Remark}
\newtheorem{remarks}[subsection]{Remarks}

\numberwithin{equation}{subsection}

% Macros
%
\def\lim{\mathop{\mathrm{lim}}\nolimits}
\def\colim{\mathop{\mathrm{colim}}\nolimits}
\def\Spec{\mathop{\mathrm{Spec}}}
\def\Hom{\mathop{\mathrm{Hom}}\nolimits}
\def\Ext{\mathop{\mathrm{Ext}}\nolimits}
\def\SheafHom{\mathop{\mathcal{H}\!\mathit{om}}\nolimits}
\def\SheafExt{\mathop{\mathcal{E}\!\mathit{xt}}\nolimits}
\def\Sch{\mathit{Sch}}
\def\Mor{\mathop{\mathrm{Mor}}\nolimits}
\def\Ob{\mathop{\mathrm{Ob}}\nolimits}
\def\Sh{\mathop{\mathit{Sh}}\nolimits}
\def\NL{\mathop{N\!L}\nolimits}
\def\CH{\mathop{\mathrm{CH}}\nolimits}
\def\proetale{{pro\text{-}\acute{e}tale}}
\def\etale{{\acute{e}tale}}
\def\QCoh{\mathit{QCoh}}
\def\Ker{\mathop{\mathrm{Ker}}}
\def\Im{\mathop{\mathrm{Im}}}
\def\Coker{\mathop{\mathrm{Coker}}}
\def\Coim{\mathop{\mathrm{Coim}}}

% Boxtimes
%
\DeclareMathSymbol{\boxtimes}{\mathbin}{AMSa}{"02}

%
% Macros for moduli stacks/spaces
%
\def\QCohstack{\mathcal{QC}\!\mathit{oh}}
\def\Cohstack{\mathcal{C}\!\mathit{oh}}
\def\Spacesstack{\mathcal{S}\!\mathit{paces}}
\def\Quotfunctor{\mathrm{Quot}}
\def\Hilbfunctor{\mathrm{Hilb}}
\def\Curvesstack{\mathcal{C}\!\mathit{urves}}
\def\Polarizedstack{\mathcal{P}\!\mathit{olarized}}
\def\Complexesstack{\mathcal{C}\!\mathit{omplexes}}
% \Pic is the operator that assigns to X its picard group, usage \Pic(X)
% \Picardstack_{X/B} denotes the Picard stack of X over B
% \Picardfunctor_{X/B} denotes the Picard functor of X over B
\def\Pic{\mathop{\mathrm{Pic}}\nolimits}
\def\Picardstack{\mathcal{P}\!\mathit{ic}}
\def\Picardfunctor{\mathrm{Pic}}
\def\Deformationcategory{\mathcal{D}\!\mathit{ef}}
\usepackage{cleveref}
\crefname{originalreference}{original reference}{original references}
\Crefname{originalreference}{Original reference}{Original references}
\makeatletter
\newlabel{cohomology-lemma-dual-perfect-complex}{{08DQ}{source}{}{sourcepointer-1}{}}
\newlabel{cohomology-lemma-dual-perfect-complex@cref}{{[originalreference][0][]08DQ}{[1][0][]source}{}{}{}}
\makeatother
\crefname{originalreference}{original reference}{original references}
\Crefname{originalreference}{Original reference}{Original references}
\makeatletter
\newlabel{cohomology-lemma-invertible-derived}{{0FPG}{source}{}{sourcepointer-2}{}}
\newlabel{cohomology-lemma-invertible-derived@cref}{{[originalreference][0][]0FPG}{[1][0][]source}{}{}{}}
\makeatother
\crefname{originalreference}{original reference}{original references}
\Crefname{originalreference}{Original reference}{Original references}
\makeatletter
\newlabel{duality-lemma-affine-flat-Noetherian-gorenstein}{{0C08}{source}{}{sourcepointer-3}{}}
\newlabel{duality-lemma-affine-flat-Noetherian-gorenstein@cref}{{[originalreference][0][]0C08}{[1][0][]source}{}{}{}}
\makeatother
\crefname{originalreference}{original reference}{original references}
\Crefname{originalreference}{Original reference}{Original references}
\makeatletter
\newlabel{duality-remark-relative-dualizing-complex}{{0B6S}{source}{}{sourcepointer-4}{}}
\newlabel{duality-remark-relative-dualizing-complex@cref}{{[originalreference][0][]0B6S}{[1][0][]source}{}{}{}}
\makeatother
\crefname{originalreference}{original reference}{original references}
\Crefname{originalreference}{Original reference}{Original references}
\makeatletter
\newlabel{duality-section-upper-shriek}{{0A9Y}{source}{}{sourcepointer-5}{}}
\newlabel{duality-section-upper-shriek@cref}{{[originalreference][0][]0A9Y}{[1][0][]source}{}{}{}}
\makeatother
\crefname{originallemma}{original lemma}{original lemmas}
\Crefname{originallemma}{Original lemma}{Original lemmas}
\makeatletter
\newlabel{lemma-compose-fourier-mukai}{{0FYS}{source}{}{sourcepointer-6}{}}
\newlabel{lemma-compose-fourier-mukai@cref}{{[originallemma][0][]0FYS}{[1][0][]source}{}{}{}}
\makeatother
\crefname{originallemma}{original lemma}{original lemmas}
\Crefname{originallemma}{Original lemma}{Original lemmas}
\makeatletter
\newlabel{lemma-fourier-mukai}{{0FYT}{source}{}{sourcepointer-7}{}}
\newlabel{lemma-fourier-mukai@cref}{{[originallemma][0][]0FYT}{[1][0][]source}{}{}{}}
\makeatother
\crefname{originallemma}{original lemma}{original lemmas}
\Crefname{originallemma}{Original lemma}{Original lemmas}
\makeatletter
\newlabel{lemma-fourier-mukai-left-adjoint}{{0FYX}{source}{}{sourcepointer-8}{}}
\newlabel{lemma-fourier-mukai-left-adjoint@cref}{{[originallemma][0][]0FYX}{[1][0][]source}{}{}{}}
\makeatother
\crefname{originallemma}{original lemma}{original lemmas}
\Crefname{originallemma}{Original lemma}{Original lemmas}
\makeatletter
\newlabel{lemma-fourier-mukai-right-adjoint}{{0FYW}{source}{}{sourcepointer-9}{}}
\newlabel{lemma-fourier-mukai-right-adjoint@cref}{{[originallemma][0][]0FYW}{[1][0][]source}{}{}{}}
\makeatother
\crefname{originalreference}{original reference}{original references}
\Crefname{originalreference}{Original reference}{Original references}
\makeatletter
\newlabel{perfect-lemma-flat-proper-perfect-direct-image}{{0B6F}{source}{}{sourcepointer-10}{}}
\newlabel{perfect-lemma-flat-proper-perfect-direct-image@cref}{{[originalreference][0][]0B6F}{[1][0][]source}{}{}{}}
\makeatother
\title{FM functor on perfect complexes, composition, and adjoint kernels}
\author{The Stacks Project Authors}
\date{stacks-project:equiv.tex:local-snapshot-20261012; GNU FDL 1.2 or later}
\begin{document}
\maketitle
\noindent\textit{Editorial note.} Parts (3)-(4) delegate the general adjunction to same-chapter lemmas carrying external citations (uncardable here); the kernel identification, base change, and invertibility arguments are in-span; part (4) is explicitly mirror-level.\par\medskip
\noindent\textit{Allowed standard prerequisites (editorial).}\begin{itemize}
\item FM functor well-definedness. lemma-fourier-mukai (same chapter, judged citation-only, not carded). Part (1) source.
\item Composition kernel. lemma-compose-fourier-mukai L26466 (carded this batch). Part (2) source and the K\textasciicircum{}prime-prime formula.
\item General adjunction kernels. lemma-fourier-mukai-right-adjoint, lemma-fourier-mukai-left-adjoint (same chapter, ext-cite filtered, not cardable). Provide the abstract adjoint kernels; identification done in-span.
\end{itemize}\medskip
\ifdefined\hypertarget\hypertarget{sourcepointer-1}{}\fi
\noindent\textit{Source reference.} Original reference 08DQ: focused-seeds/geometry/S9-eqv-fourier-mukai-flat-proper-over-noetherian/ref-tags.txt, source lines 5--5. Navigation pointer only: the counted proof cites this statement by \textbackslash{}ref as a standing input; the deduction the counted proof has to make is carried out inside the counted span, and no part of that deduction is delegated to this pointer.\par\medskip
\ifdefined\hypertarget\hypertarget{sourcepointer-2}{}\fi
\noindent\textit{Source reference.} Original reference 0FPG: focused-seeds/geometry/S9-eqv-fourier-mukai-flat-proper-over-noetherian/ref-tags.txt, source lines 6--6. Navigation pointer only: the counted proof cites this statement by \textbackslash{}ref as a standing input; the deduction the counted proof has to make is carried out inside the counted span, and no part of that deduction is delegated to this pointer.\par\medskip
\ifdefined\hypertarget\hypertarget{sourcepointer-3}{}\fi
\noindent\textit{Source reference.} Original reference 0C08: focused-seeds/geometry/S9-eqv-fourier-mukai-flat-proper-over-noetherian/ref-tags.txt, source lines 7--7. Navigation pointer only: the counted proof cites this statement by \textbackslash{}ref as a standing input; the deduction the counted proof has to make is carried out inside the counted span, and no part of that deduction is delegated to this pointer.\par\medskip
\ifdefined\hypertarget\hypertarget{sourcepointer-4}{}\fi
\noindent\textit{Source reference.} Original reference 0B6S: focused-seeds/geometry/S9-eqv-fourier-mukai-flat-proper-over-noetherian/ref-tags.txt, source lines 8--8. Navigation pointer only: the counted proof cites this statement by \textbackslash{}ref as a standing input; the deduction the counted proof has to make is carried out inside the counted span, and no part of that deduction is delegated to this pointer.\par\medskip
\ifdefined\hypertarget\hypertarget{sourcepointer-5}{}\fi
\noindent\textit{Source reference.} Original reference 0A9Y: focused-seeds/geometry/S9-eqv-fourier-mukai-flat-proper-over-noetherian/ref-tags.txt, source lines 9--9. Navigation pointer only: the counted proof cites this statement by \textbackslash{}ref as a standing input; the deduction the counted proof has to make is carried out inside the counted span, and no part of that deduction is delegated to this pointer.\par\medskip
\ifdefined\hypertarget\hypertarget{sourcepointer-6}{}\fi
\noindent\textit{Source reference.} Original lemma 0FYS: sources/stacks-equiv/pinned-original-source.tex, source lines 1200--1200. Navigation pointer only: the counted proof cites this statement by \textbackslash{}ref as a standing input; the deduction the counted proof has to make is carried out inside the counted span, and no part of that deduction is delegated to this pointer.\par\medskip
\ifdefined\hypertarget\hypertarget{sourcepointer-7}{}\fi
\noindent\textit{Source reference.} Original lemma 0FYT: sources/stacks-equiv/pinned-original-source.tex, source lines 1282--1282. Navigation pointer only: the counted proof cites this statement by \textbackslash{}ref as a standing input; the deduction the counted proof has to make is carried out inside the counted span, and no part of that deduction is delegated to this pointer.\par\medskip
\ifdefined\hypertarget\hypertarget{sourcepointer-8}{}\fi
\noindent\textit{Source reference.} Original lemma 0FYX: sources/stacks-equiv/pinned-original-source.tex, source lines 1577--1577. Navigation pointer only: the counted proof cites this statement by \textbackslash{}ref as a standing input; the deduction the counted proof has to make is carried out inside the counted span, and no part of that deduction is delegated to this pointer.\par\medskip
\ifdefined\hypertarget\hypertarget{sourcepointer-9}{}\fi
\noindent\textit{Source reference.} Original lemma 0FYW: sources/stacks-equiv/pinned-original-source.tex, source lines 1487--1487. Navigation pointer only: the counted proof cites this statement by \textbackslash{}ref as a standing input; the deduction the counted proof has to make is carried out inside the counted span, and no part of that deduction is delegated to this pointer.\par\medskip
\ifdefined\hypertarget\hypertarget{sourcepointer-10}{}\fi
\noindent\textit{Source reference.} Original reference 0B6F: focused-seeds/geometry/S9-eqv-fourier-mukai-flat-proper-over-noetherian/ref-tags.txt, source lines 10--10. Navigation pointer only: the counted proof cites this statement by \textbackslash{}ref as a standing input; the deduction the counted proof has to make is carried out inside the counted span, and no part of that deduction is delegated to this pointer.\par\medskip
\begin{lemma}
\label{lemma-fourier-mukai-flat-proper-over-noetherian}
Let $S$ be a Noetherian scheme.
\begin{enumerate}
\item For $X$, $Y$ proper and flat over $S$ and $K$ in
$D_{perf}(\mathcal{O}_{X \times_S Y})$ we obtain a Fourier-Mukai functor
$\Phi_K : D_{perf}(\mathcal{O}_X) \to D_{perf}(\mathcal{O}_Y)$.
\item For $X$, $Y$, $Z$ proper and flat over $S$, $K \in
D_{perf}(\mathcal{O}_{X \times_S Y})$, $K' \in
D_{perf}(\mathcal{O}_{Y \times_S Z})$ the composition
$\Phi_{K'} \circ \Phi_K : D_{perf}(\mathcal{O}_X) \to D_{perf}(\mathcal{O}_Z)$
is equal to $\Phi_{K''}$ with $K'' \in D_{perf}(\mathcal{O}_{X \times_S Z})$
computed as in Lemma \ref{lemma-compose-fourier-mukai},
\item For $X$, $Y$, $K$, $\Phi_K$ as in (1) if $X \to S$ is Gorenstein, then
$\Phi_{K'} : D_{perf}(\mathcal{O}_Y) \to D_{perf}(\mathcal{O}_X)$ is a right
adjoint to $\Phi_K$ where $K' \in D_{perf}(\mathcal{O}_{Y \times_S X})$
is the pullback of $L\text{pr}_1^*\omega_{X/S}^\bullet
\otimes_{\mathcal{O}_{X \times_S Y}}^\mathbf{L} K^\vee$ by
$Y \times_S X \to X \times_S Y$.
\item For $X$, $Y$, $K$, $\Phi_K$ as in (1) if $Y \to S$ is Gorenstein, then
$\Phi_{K''} : D_{perf}(\mathcal{O}_Y) \to D_{perf}(\mathcal{O}_X)$ is a left
adjoint to $\Phi_K$ where $K'' \in D_{perf}(\mathcal{O}_{Y \times_S X})$
is the pullback of $L\text{pr}_2^*\omega_{Y/S}^\bullet
\otimes_{\mathcal{O}_{X \times_S Y}}^\mathbf{L} K^\vee$ by
$Y \times_S X \to X \times_S Y$.
\end{enumerate}
\end{lemma}

\begin{proof}
Part (1) is immediate from Lemma \ref{lemma-fourier-mukai} part (4).

\medskip\noindent
Part (2) follows from Lemma \ref{lemma-compose-fourier-mukai} and the
fact that
$K'' = R\text{pr}_{13, *}(
L\text{pr}_{12}^*K
\otimes_{\mathcal{O}_{X \times_S Y \times_S Z}}^\mathbf{L}
L\text{pr}_{23}^*K')$ is perfect for example by
Derived Categories of Schemes, Lemma
\ref{perfect-lemma-flat-proper-perfect-direct-image}.

\medskip\noindent
The adjointness in part (3) on all complexes with quasi-coherent cohomology
sheaves follows from Lemma \ref{lemma-fourier-mukai-right-adjoint} with
$K'$ equal to the pullback of
$R\SheafHom_{\mathcal{O}_{X \times_S Y}}(K, a(\mathcal{O}_Y))$
by $Y \times_S X \to X \times_S Y$ where $a$ is the right adjoint
to $R\text{pr}_{2, *} : D_\QCoh(\mathcal{O}_{X \times_S Y}) \to
D_\QCoh(\mathcal{O}_Y)$. Denote $f : X \to S$ the structure morphism of $X$.
Since $f$ is proper the functor
$f^! : D_\QCoh^+(\mathcal{O}_S) \to D_\QCoh^+(\mathcal{O}_X)$
is the restriction to $D_\QCoh^+(\mathcal{O}_S)$
of the right adjoint to
$Rf_* : D_\QCoh(\mathcal{O}_X) \to D_\QCoh(\mathcal{O}_S)$, see
Duality for Schemes, Section \ref{duality-section-upper-shriek}.
Hence the relative dualizing complex $\omega_{X/S}^\bullet$ as defined in
Duality for Schemes, Remark
\ref{duality-remark-relative-dualizing-complex}
is equal to $\omega_{X/S}^\bullet = f^!\mathcal{O}_S$.
Since formation of the relative dualizing complex
commutes with base change (see Duality for Schemes, Remark
\ref{duality-remark-relative-dualizing-complex}) we see that
$a(\mathcal{O}_Y) = L\text{pr}_1^*\omega_{X/S}^\bullet$.
Thus
$$
R\SheafHom_{\mathcal{O}_{X \times_S Y}}(K, a(\mathcal{O}_Y))
\cong
L\text{pr}_1^*\omega_{X/S}^\bullet
\otimes_{\mathcal{O}_{X \times_S Y}}^\mathbf{L} K^\vee
$$
by Cohomology, Lemma \ref{cohomology-lemma-dual-perfect-complex}.
Finally, since $X \to S$ is assumed Gorenstein the relative dualizing complex
is invertible: this follows from Duality for Schemes, Lemma
\ref{duality-lemma-affine-flat-Noetherian-gorenstein}.
We conclude that $\omega_{X/S}^\bullet$ is perfect
(Cohomology, Lemma \ref{cohomology-lemma-invertible-derived})
and hence $K'$ is perfect.
Therefore $\Phi_{K'}$ does indeed map $D_{perf}(\mathcal{O}_Y)$
into $D_{perf}(\mathcal{O}_X)$ which finishes the proof of (3).

\medskip\noindent
The proof of (4) is the same as the proof of (3) except one uses
Lemma \ref{lemma-fourier-mukai-left-adjoint} instead of
Lemma \ref{lemma-fourier-mukai-right-adjoint}.
\end{proof}

\end{document}
